
%!TEX root = ../csuthesis_main.tex
% 设置英文摘要
\keywordsen{Quantum Machine Learning; Audio Classification; Quantum Time-Series Encoding; Quantum Time-Frequency Analysis Neural Network; Knowledge Distillation}
\categoryen{TP183}
%\itemcountcn{There are \totalfigures\ figures, \totaltables\ tables, and \total{citnum}\ citations in this thesis.}
%\addcontentsline{toc}{section}{ABSTRACT}
\begin{abstracten}\setlength{\baselineskip}{20pt}%\renewcommand{\baselinestretch}{1.0}
	Audio classification is an important research direction in the fields of artificial intelligence and digital signal processing.
	In recent years, classical deep learning models have made significant progress but also face challenges regarding massive parameter sizes and computational bottlenecks.
	With the development of quantum computing, quantum audio processing has emerged, leveraging quantum superposition and entanglement to achieve exponentially accelerated feature extraction at a markedly reduced parameter scale.
	However, existing methods often ignore the temporal evolution characteristics of audio signals, and face difficulties in training deep circuits on Noisy Intermediate-Scale Quantum (NISQ) devices.
	Therefore, exploring a quantum neural network architecture that can efficiently represent temporal features and possess good convergence has become a crucial research topic.
	This thesis proposes an audio classification and optimization framework for quantum neural networks based on time-series encoding.
	It systematically explores two main dimensions: quantum feature extraction network construction and cross-paradigm training strategies.
	The main research contributions of this thesis are as follows:

(1) A highly parameter-efficient Quantum Time-Frequency Analysis Neural Network model based on Quantum Time-Series Encoding is proposed. To address the limitations of traditional encoding in capturing dynamic temporal features, this method designs a quantum time-series encoding scheme based on dual-qubit sequence entanglement. It globally entangles discretized Mel-Frequency Cepstral Coefficients (MFCCs) with audio time indices, effectively solving the challenge of mapping temporal features into a high-dimensional quantum space. Building upon this encoding, an end-to-end architecture comprising a quantum convolutional layer, a quantum pooling layer, and a quantum measurement module is constructed by integrating the Hardware-Efficient Ansatz. Experimental results on the GTZAN dataset demonstrate that the model achieves a classification accuracy of 75\% with merely 33 fixed trainable parameters, exhibiting significant advantages in parameter efficiency and feature capture capability.
	
	(2) A cross-paradigm knowledge distillation strategy, structured around a ``classical teacher-quantum student'' framework, is introduced for optimization. 
    Variational quantum circuits operating on NISQ devices frequently encounter local optima and barren plateaus. The proposed strategy employs a pre-trained classical Convolutional Neural Network as a teacher network, which extracts high-dimensional acoustic features as continuous probability distributions. 
    A temperature is applied to formulate a joint loss function combining hard cross-entropy with soft KL divergence. This function directs the parameter updates within the highly compact quantum circuit. 
    Such a soft-label mechanism serves to smooth the loss landscape throughout the quantum evolution process. 
    Subsequent testing confirms that this distillation approach successfully narrows the generalization gap inherent in the pure quantum model. 
    Specifically, the test set accuracy increases to 76.17\%, accompanied by an 82.67\% recall rate and an AUC of 0.8319. A noticeable reduction in the Brier score is also observed.
	Overall, these findings validate the strategy's capability to accelerate convergence, ensure reliable predictions, and bolster the noise robustness of the quantum circuit.
	% 【修改终点】
\end{abstracten}